% paper35-partiality-models.tex — Partiality Model Reconstruction:
%   Handling Non-Termination and Exceptions
% Paper 35 of the Judgment Geometry series.
% Compile: pdflatex -interaction=nonstopmode paper35-partiality-models.tex

\documentclass[11pt]{article}
\usepackage{lmodern}
% jugeo-common.tex — shared preamble for all Judgment Geometry papers
% Include via: % jugeo-common.tex — shared preamble for all Judgment Geometry papers
% Include via: % jugeo-common.tex — shared preamble for all Judgment Geometry papers
% Include via: \input{jugeo-common}

% ── Packages ────────────────────────────────────────────────────────────
\usepackage{amsmath,amssymb,amsthm,mathtools}
\usepackage{tikz-cd}
\usepackage{listings}
\usepackage{xcolor}
\usepackage{xspace}
\usepackage{booktabs}
\usepackage{multirow}
\usepackage{enumitem}
\usepackage[expansion=false]{microtype}
\usepackage{hyperref}
\usepackage{cleveref}
% \usepackage{thmtools}  % disabled: not available in this TeX installation
\usepackage{float}
\usepackage{caption}
\usepackage{subcaption}
\usepackage{tabularx}
\usepackage{stmaryrd}       % \llbracket, \rrbracket
\usepackage{bbm}            % \mathbbm{1}

% ── Page geometry ───────────────────────────────────────────────────────
\usepackage[margin=1in]{geometry}

% ── Theorem environments ────────────────────────────────────────────────
\theoremstyle{plain}
\newtheorem{theorem}{Theorem}[section]
\newtheorem{lemma}[theorem]{Lemma}
\newtheorem{proposition}[theorem]{Proposition}
\newtheorem{corollary}[theorem]{Corollary}

\theoremstyle{definition}
\newtheorem{definition}[theorem]{Definition}
\newtheorem{example}[theorem]{Example}
\newtheorem{construction}[theorem]{Construction}

\theoremstyle{remark}
\newtheorem{remark}[theorem]{Remark}
\newtheorem{notation}[theorem]{Notation}

% ── Python listings style ──────────────────────────────────────────────
\definecolor{jugeoBlue}{HTML}{2563EB}
\definecolor{jugeoGreen}{HTML}{059669}
\definecolor{jugeoRed}{HTML}{DC2626}
\definecolor{jugeoPurple}{HTML}{7C3AED}
\definecolor{jugeoGray}{HTML}{6B7280}
\definecolor{jugeoBg}{HTML}{F8FAFC}

\lstdefinestyle{jugeo-python}{
  language=Python,
  basicstyle=\ttfamily\small,
  keywordstyle=\color{jugeoBlue}\bfseries,
  stringstyle=\color{jugeoGreen},
  commentstyle=\color{jugeoGray}\itshape,
  numberstyle=\tiny\color{jugeoGray},
  backgroundcolor=\color{jugeoBg},
  numbers=left,
  numbersep=6pt,
  frame=single,
  framerule=0.4pt,
  rulecolor=\color{jugeoGray!40},
  breaklines=true,
  breakatwhitespace=true,
  showstringspaces=false,
  tabsize=4,
  xleftmargin=1.5em,
  framexleftmargin=1.5em,
  morekeywords={dataclass,frozen,field,Enum,IntEnum,auto,Any,
                Mapping,Sequence,Optional,match,case},
  literate={→}{{$\to$}}1 {←}{{$\leftarrow$}}1
           {≤}{{$\leq$}}1 {≥}{{$\geq$}}1 {≠}{{$\neq$}}1
           {∈}{{$\in$}}1 {∀}{{$\forall$}}1 {∃}{{$\exists$}}1
           {⊕}{{$\oplus$}}1 {⊖}{{$\ominus$}}1,
  escapeinside={(*@}{@*)},
}
\lstset{style=jugeo-python}

\lstdefinestyle{jugeo-output}{
  basicstyle=\ttfamily\small,
  backgroundcolor=\color{jugeoBg},
  frame=single,
  framerule=0.4pt,
  rulecolor=\color{jugeoGray!40},
  breaklines=true,
  showstringspaces=false,
  xleftmargin=1.5em,
  framexleftmargin=0.5em,
  numbers=none,
}

% ── JuGeo-specific macros ──────────────────────────────────────────────

% Judgment 8-tuple
\newcommand{\J}{\mathcal{J}}
\newcommand{\Jud}[8]{(#1,\,#2,\,#3,\,#4,\,#5,\,#6,\,#7,\,#8)}
\newcommand{\JudShort}[1]{\J_{#1}}

% Components
\newcommand{\coord}{\mathsf{c}}            % coordinate
\newcommand{\prop}{\varphi}                 % proposition
\newcommand{\carrier}{\mathcal{A}}          % carrier type
\newcommand{\evid}{\mathcal{E}}             % evidence bundle
\newcommand{\oblig}{\mathcal{O}}            % obligations
\newcommand{\obstr}{\mathcal{B}}            % obstructions
\newcommand{\trust}{\mathcal{T}}            % trust annotation
\newcommand{\prov}{\Pi}                     % provenance

% Trust algebra
\newcommand{\Talg}{\mathfrak{T}}            % trust ordered algebra
\newcommand{\Eadm}{\mathcal{E}_{\mathrm{adm}}}
\newcommand{\tleq}{\preceq}                 % trust ordering
\newcommand{\tjoin}{\oplus}                 % conservative join
\newcommand{\tatten}{\ominus}               % attenuation
\newcommand{\tpromote}{\uparrow_\pi}        % promotion
\newcommand{\tdemote}{\downarrow_\chi}       % demotion

% Trust tiers
\newcommand{\Tcontra}{\ensuremath{\mathsf{CONTRADICTED}}}
\newcommand{\Tunver}{\ensuremath{\mathsf{UNVERIFIED}}}
\newcommand{\Tcopilot}{\ensuremath{\mathsf{COPILOT}}}
\newcommand{\Toracle}{\ensuremath{\mathsf{ORACLE}}}
\newcommand{\Truntime}{\ensuremath{\mathsf{RUNTIME}}}
\newcommand{\Tsolver}{\ensuremath{\mathsf{SOLVER}}}
\newcommand{\Tproof}{\ensuremath{\mathsf{PROOF}}}

% Site and geometry
\newcommand{\Site}{\mathcal{S}}             % semantic site
\newcommand{\Cov}{\mathrm{Cov}}            % covering family
\newcommand{\Ob}{\mathrm{Ob}}              % objects
\newcommand{\Mor}{\mathrm{Mor}}            % morphisms
\newcommand{\res}{\mathrm{res}}            % restriction
\newcommand{\Cech}{\check{C}}              % Čech complex
\newcommand{\Hcoh}[1]{H^{#1}}             % cohomology group

% Descent
\newcommand{\descent}{\mathrm{desc}}
\newcommand{\glue}{\mathrm{glue}}
\newcommand{\overlap}[2]{#1 \cap #2}

% Semantic moves
\newcommand{\VERIFY}{\textsc{Verify}}
\newcommand{\CONSTRUCT}{\textsc{Construct}}
\newcommand{\NORMALIZE}{\textsc{Normalize}}
\newcommand{\GLUE}{\textsc{Glue}}
\newcommand{\RESTRICT}{\textsc{Restrict}}
\newcommand{\TRANSPORT}{\textsc{Transport}}
\newcommand{\REPAIR}{\textsc{Repair}}
\newcommand{\PROMOTE}{\textsc{Promote}}
\newcommand{\CHALLENGE}{\textsc{Challenge}}
\newcommand{\ESCALATE}{\textsc{Escalate}}

% Fragments
\newcommand{\QFLIA}{\texttt{QF\_LIA}}
\newcommand{\QFLRA}{\texttt{QF\_LRA}}
\newcommand{\QFBV}{\texttt{QF\_BV}}
\newcommand{\QFUF}{\texttt{QF\_UF}}

% Misc
\newcommand{\Python}{\textsc{Python}}
\newcommand{\JuGeo}{\textsc{JuGeo}}
\newcommand{\LEAN}{\textsc{Lean}\,4}
\newcommand{\Fstar}{F$^{\star}$}
\newcommand{\Coq}{\textsc{Coq}}
\newcommand{\Dafny}{\textsc{Dafny}}

% Common shortcuts
\newcommand{\ie}{i.e.\@\xspace}
\newcommand{\eg}{e.g.\@\xspace}
\newcommand{\cf}{cf.\@\xspace}
\newcommand{\wrt}{w.r.t.\@\xspace}

% Evaluation shortcuts
\newcommand{\numcases}{300}
\newcommand{\numspec}{100}
\newcommand{\numeq}{100}
\newcommand{\numbug}{100}
\newcommand{\medtime}{6.5\,ms}
\newcommand{\totaltime}{1.949\,s}


% ── Packages ────────────────────────────────────────────────────────────
\usepackage{amsmath,amssymb,amsthm,mathtools}
\usepackage{tikz-cd}
\usepackage{listings}
\usepackage{xcolor}
\usepackage{xspace}
\usepackage{booktabs}
\usepackage{multirow}
\usepackage{enumitem}
\usepackage[expansion=false]{microtype}
\usepackage{hyperref}
\usepackage{cleveref}
% \usepackage{thmtools}  % disabled: not available in this TeX installation
\usepackage{float}
\usepackage{caption}
\usepackage{subcaption}
\usepackage{tabularx}
\usepackage{stmaryrd}       % \llbracket, \rrbracket
\usepackage{bbm}            % \mathbbm{1}

% ── Page geometry ───────────────────────────────────────────────────────
\usepackage[margin=1in]{geometry}

% ── Theorem environments ────────────────────────────────────────────────
\theoremstyle{plain}
\newtheorem{theorem}{Theorem}[section]
\newtheorem{lemma}[theorem]{Lemma}
\newtheorem{proposition}[theorem]{Proposition}
\newtheorem{corollary}[theorem]{Corollary}

\theoremstyle{definition}
\newtheorem{definition}[theorem]{Definition}
\newtheorem{example}[theorem]{Example}
\newtheorem{construction}[theorem]{Construction}

\theoremstyle{remark}
\newtheorem{remark}[theorem]{Remark}
\newtheorem{notation}[theorem]{Notation}

% ── Python listings style ──────────────────────────────────────────────
\definecolor{jugeoBlue}{HTML}{2563EB}
\definecolor{jugeoGreen}{HTML}{059669}
\definecolor{jugeoRed}{HTML}{DC2626}
\definecolor{jugeoPurple}{HTML}{7C3AED}
\definecolor{jugeoGray}{HTML}{6B7280}
\definecolor{jugeoBg}{HTML}{F8FAFC}

\lstdefinestyle{jugeo-python}{
  language=Python,
  basicstyle=\ttfamily\small,
  keywordstyle=\color{jugeoBlue}\bfseries,
  stringstyle=\color{jugeoGreen},
  commentstyle=\color{jugeoGray}\itshape,
  numberstyle=\tiny\color{jugeoGray},
  backgroundcolor=\color{jugeoBg},
  numbers=left,
  numbersep=6pt,
  frame=single,
  framerule=0.4pt,
  rulecolor=\color{jugeoGray!40},
  breaklines=true,
  breakatwhitespace=true,
  showstringspaces=false,
  tabsize=4,
  xleftmargin=1.5em,
  framexleftmargin=1.5em,
  morekeywords={dataclass,frozen,field,Enum,IntEnum,auto,Any,
                Mapping,Sequence,Optional,match,case},
  literate={→}{{$\to$}}1 {←}{{$\leftarrow$}}1
           {≤}{{$\leq$}}1 {≥}{{$\geq$}}1 {≠}{{$\neq$}}1
           {∈}{{$\in$}}1 {∀}{{$\forall$}}1 {∃}{{$\exists$}}1
           {⊕}{{$\oplus$}}1 {⊖}{{$\ominus$}}1,
  escapeinside={(*@}{@*)},
}
\lstset{style=jugeo-python}

\lstdefinestyle{jugeo-output}{
  basicstyle=\ttfamily\small,
  backgroundcolor=\color{jugeoBg},
  frame=single,
  framerule=0.4pt,
  rulecolor=\color{jugeoGray!40},
  breaklines=true,
  showstringspaces=false,
  xleftmargin=1.5em,
  framexleftmargin=0.5em,
  numbers=none,
}

% ── JuGeo-specific macros ──────────────────────────────────────────────

% Judgment 8-tuple
\newcommand{\J}{\mathcal{J}}
\newcommand{\Jud}[8]{(#1,\,#2,\,#3,\,#4,\,#5,\,#6,\,#7,\,#8)}
\newcommand{\JudShort}[1]{\J_{#1}}

% Components
\newcommand{\coord}{\mathsf{c}}            % coordinate
\newcommand{\prop}{\varphi}                 % proposition
\newcommand{\carrier}{\mathcal{A}}          % carrier type
\newcommand{\evid}{\mathcal{E}}             % evidence bundle
\newcommand{\oblig}{\mathcal{O}}            % obligations
\newcommand{\obstr}{\mathcal{B}}            % obstructions
\newcommand{\trust}{\mathcal{T}}            % trust annotation
\newcommand{\prov}{\Pi}                     % provenance

% Trust algebra
\newcommand{\Talg}{\mathfrak{T}}            % trust ordered algebra
\newcommand{\Eadm}{\mathcal{E}_{\mathrm{adm}}}
\newcommand{\tleq}{\preceq}                 % trust ordering
\newcommand{\tjoin}{\oplus}                 % conservative join
\newcommand{\tatten}{\ominus}               % attenuation
\newcommand{\tpromote}{\uparrow_\pi}        % promotion
\newcommand{\tdemote}{\downarrow_\chi}       % demotion

% Trust tiers
\newcommand{\Tcontra}{\ensuremath{\mathsf{CONTRADICTED}}}
\newcommand{\Tunver}{\ensuremath{\mathsf{UNVERIFIED}}}
\newcommand{\Tcopilot}{\ensuremath{\mathsf{COPILOT}}}
\newcommand{\Toracle}{\ensuremath{\mathsf{ORACLE}}}
\newcommand{\Truntime}{\ensuremath{\mathsf{RUNTIME}}}
\newcommand{\Tsolver}{\ensuremath{\mathsf{SOLVER}}}
\newcommand{\Tproof}{\ensuremath{\mathsf{PROOF}}}

% Site and geometry
\newcommand{\Site}{\mathcal{S}}             % semantic site
\newcommand{\Cov}{\mathrm{Cov}}            % covering family
\newcommand{\Ob}{\mathrm{Ob}}              % objects
\newcommand{\Mor}{\mathrm{Mor}}            % morphisms
\newcommand{\res}{\mathrm{res}}            % restriction
\newcommand{\Cech}{\check{C}}              % Čech complex
\newcommand{\Hcoh}[1]{H^{#1}}             % cohomology group

% Descent
\newcommand{\descent}{\mathrm{desc}}
\newcommand{\glue}{\mathrm{glue}}
\newcommand{\overlap}[2]{#1 \cap #2}

% Semantic moves
\newcommand{\VERIFY}{\textsc{Verify}}
\newcommand{\CONSTRUCT}{\textsc{Construct}}
\newcommand{\NORMALIZE}{\textsc{Normalize}}
\newcommand{\GLUE}{\textsc{Glue}}
\newcommand{\RESTRICT}{\textsc{Restrict}}
\newcommand{\TRANSPORT}{\textsc{Transport}}
\newcommand{\REPAIR}{\textsc{Repair}}
\newcommand{\PROMOTE}{\textsc{Promote}}
\newcommand{\CHALLENGE}{\textsc{Challenge}}
\newcommand{\ESCALATE}{\textsc{Escalate}}

% Fragments
\newcommand{\QFLIA}{\texttt{QF\_LIA}}
\newcommand{\QFLRA}{\texttt{QF\_LRA}}
\newcommand{\QFBV}{\texttt{QF\_BV}}
\newcommand{\QFUF}{\texttt{QF\_UF}}

% Misc
\newcommand{\Python}{\textsc{Python}}
\newcommand{\JuGeo}{\textsc{JuGeo}}
\newcommand{\LEAN}{\textsc{Lean}\,4}
\newcommand{\Fstar}{F$^{\star}$}
\newcommand{\Coq}{\textsc{Coq}}
\newcommand{\Dafny}{\textsc{Dafny}}

% Common shortcuts
\newcommand{\ie}{i.e.\@\xspace}
\newcommand{\eg}{e.g.\@\xspace}
\newcommand{\cf}{cf.\@\xspace}
\newcommand{\wrt}{w.r.t.\@\xspace}

% Evaluation shortcuts
\newcommand{\numcases}{300}
\newcommand{\numspec}{100}
\newcommand{\numeq}{100}
\newcommand{\numbug}{100}
\newcommand{\medtime}{6.5\,ms}
\newcommand{\totaltime}{1.949\,s}


% ── Packages ────────────────────────────────────────────────────────────
\usepackage{amsmath,amssymb,amsthm,mathtools}
\usepackage{tikz-cd}
\usepackage{listings}
\usepackage{xcolor}
\usepackage{xspace}
\usepackage{booktabs}
\usepackage{multirow}
\usepackage{enumitem}
\usepackage[expansion=false]{microtype}
\usepackage{hyperref}
\usepackage{cleveref}
% \usepackage{thmtools}  % disabled: not available in this TeX installation
\usepackage{float}
\usepackage{caption}
\usepackage{subcaption}
\usepackage{tabularx}
\usepackage{stmaryrd}       % \llbracket, \rrbracket
\usepackage{bbm}            % \mathbbm{1}

% ── Page geometry ───────────────────────────────────────────────────────
\usepackage[margin=1in]{geometry}

% ── Theorem environments ────────────────────────────────────────────────
\theoremstyle{plain}
\newtheorem{theorem}{Theorem}[section]
\newtheorem{lemma}[theorem]{Lemma}
\newtheorem{proposition}[theorem]{Proposition}
\newtheorem{corollary}[theorem]{Corollary}

\theoremstyle{definition}
\newtheorem{definition}[theorem]{Definition}
\newtheorem{example}[theorem]{Example}
\newtheorem{construction}[theorem]{Construction}

\theoremstyle{remark}
\newtheorem{remark}[theorem]{Remark}
\newtheorem{notation}[theorem]{Notation}

% ── Python listings style ──────────────────────────────────────────────
\definecolor{jugeoBlue}{HTML}{2563EB}
\definecolor{jugeoGreen}{HTML}{059669}
\definecolor{jugeoRed}{HTML}{DC2626}
\definecolor{jugeoPurple}{HTML}{7C3AED}
\definecolor{jugeoGray}{HTML}{6B7280}
\definecolor{jugeoBg}{HTML}{F8FAFC}

\lstdefinestyle{jugeo-python}{
  language=Python,
  basicstyle=\ttfamily\small,
  keywordstyle=\color{jugeoBlue}\bfseries,
  stringstyle=\color{jugeoGreen},
  commentstyle=\color{jugeoGray}\itshape,
  numberstyle=\tiny\color{jugeoGray},
  backgroundcolor=\color{jugeoBg},
  numbers=left,
  numbersep=6pt,
  frame=single,
  framerule=0.4pt,
  rulecolor=\color{jugeoGray!40},
  breaklines=true,
  breakatwhitespace=true,
  showstringspaces=false,
  tabsize=4,
  xleftmargin=1.5em,
  framexleftmargin=1.5em,
  morekeywords={dataclass,frozen,field,Enum,IntEnum,auto,Any,
                Mapping,Sequence,Optional,match,case},
  literate={→}{{$\to$}}1 {←}{{$\leftarrow$}}1
           {≤}{{$\leq$}}1 {≥}{{$\geq$}}1 {≠}{{$\neq$}}1
           {∈}{{$\in$}}1 {∀}{{$\forall$}}1 {∃}{{$\exists$}}1
           {⊕}{{$\oplus$}}1 {⊖}{{$\ominus$}}1,
  escapeinside={(*@}{@*)},
}
\lstset{style=jugeo-python}

\lstdefinestyle{jugeo-output}{
  basicstyle=\ttfamily\small,
  backgroundcolor=\color{jugeoBg},
  frame=single,
  framerule=0.4pt,
  rulecolor=\color{jugeoGray!40},
  breaklines=true,
  showstringspaces=false,
  xleftmargin=1.5em,
  framexleftmargin=0.5em,
  numbers=none,
}

% ── JuGeo-specific macros ──────────────────────────────────────────────

% Judgment 8-tuple
\newcommand{\J}{\mathcal{J}}
\newcommand{\Jud}[8]{(#1,\,#2,\,#3,\,#4,\,#5,\,#6,\,#7,\,#8)}
\newcommand{\JudShort}[1]{\J_{#1}}

% Components
\newcommand{\coord}{\mathsf{c}}            % coordinate
\newcommand{\prop}{\varphi}                 % proposition
\newcommand{\carrier}{\mathcal{A}}          % carrier type
\newcommand{\evid}{\mathcal{E}}             % evidence bundle
\newcommand{\oblig}{\mathcal{O}}            % obligations
\newcommand{\obstr}{\mathcal{B}}            % obstructions
\newcommand{\trust}{\mathcal{T}}            % trust annotation
\newcommand{\prov}{\Pi}                     % provenance

% Trust algebra
\newcommand{\Talg}{\mathfrak{T}}            % trust ordered algebra
\newcommand{\Eadm}{\mathcal{E}_{\mathrm{adm}}}
\newcommand{\tleq}{\preceq}                 % trust ordering
\newcommand{\tjoin}{\oplus}                 % conservative join
\newcommand{\tatten}{\ominus}               % attenuation
\newcommand{\tpromote}{\uparrow_\pi}        % promotion
\newcommand{\tdemote}{\downarrow_\chi}       % demotion

% Trust tiers
\newcommand{\Tcontra}{\ensuremath{\mathsf{CONTRADICTED}}}
\newcommand{\Tunver}{\ensuremath{\mathsf{UNVERIFIED}}}
\newcommand{\Tcopilot}{\ensuremath{\mathsf{COPILOT}}}
\newcommand{\Toracle}{\ensuremath{\mathsf{ORACLE}}}
\newcommand{\Truntime}{\ensuremath{\mathsf{RUNTIME}}}
\newcommand{\Tsolver}{\ensuremath{\mathsf{SOLVER}}}
\newcommand{\Tproof}{\ensuremath{\mathsf{PROOF}}}

% Site and geometry
\newcommand{\Site}{\mathcal{S}}             % semantic site
\newcommand{\Cov}{\mathrm{Cov}}            % covering family
\newcommand{\Ob}{\mathrm{Ob}}              % objects
\newcommand{\Mor}{\mathrm{Mor}}            % morphisms
\newcommand{\res}{\mathrm{res}}            % restriction
\newcommand{\Cech}{\check{C}}              % Čech complex
\newcommand{\Hcoh}[1]{H^{#1}}             % cohomology group

% Descent
\newcommand{\descent}{\mathrm{desc}}
\newcommand{\glue}{\mathrm{glue}}
\newcommand{\overlap}[2]{#1 \cap #2}

% Semantic moves
\newcommand{\VERIFY}{\textsc{Verify}}
\newcommand{\CONSTRUCT}{\textsc{Construct}}
\newcommand{\NORMALIZE}{\textsc{Normalize}}
\newcommand{\GLUE}{\textsc{Glue}}
\newcommand{\RESTRICT}{\textsc{Restrict}}
\newcommand{\TRANSPORT}{\textsc{Transport}}
\newcommand{\REPAIR}{\textsc{Repair}}
\newcommand{\PROMOTE}{\textsc{Promote}}
\newcommand{\CHALLENGE}{\textsc{Challenge}}
\newcommand{\ESCALATE}{\textsc{Escalate}}

% Fragments
\newcommand{\QFLIA}{\texttt{QF\_LIA}}
\newcommand{\QFLRA}{\texttt{QF\_LRA}}
\newcommand{\QFBV}{\texttt{QF\_BV}}
\newcommand{\QFUF}{\texttt{QF\_UF}}

% Misc
\newcommand{\Python}{\textsc{Python}}
\newcommand{\JuGeo}{\textsc{JuGeo}}
\newcommand{\LEAN}{\textsc{Lean}\,4}
\newcommand{\Fstar}{F$^{\star}$}
\newcommand{\Coq}{\textsc{Coq}}
\newcommand{\Dafny}{\textsc{Dafny}}

% Common shortcuts
\newcommand{\ie}{i.e.\@\xspace}
\newcommand{\eg}{e.g.\@\xspace}
\newcommand{\cf}{cf.\@\xspace}
\newcommand{\wrt}{w.r.t.\@\xspace}

% Evaluation shortcuts
\newcommand{\numcases}{300}
\newcommand{\numspec}{100}
\newcommand{\numeq}{100}
\newcommand{\numbug}{100}
\newcommand{\medtime}{6.5\,ms}
\newcommand{\totaltime}{1.949\,s}

% experiment-data.tex — AUTO-GENERATED from real jugeo CLI runs
% DO NOT EDIT — regenerate with: python3 experiments/run_universal_metrics.py
% Generated from 30 programs across 6 domains

\newcommand{\expTotalPrograms}{30}
\newcommand{\expVerified}{30}
\newcommand{\expAccuracy}{100.0\%}
\newcommand{\expCoordsMin}{2}
\newcommand{\expCoordsMax}{10}
\newcommand{\expCoordsMean}{5.2}
\newcommand{\expCoordsMedian}{5.5}
\newcommand{\expPropsMin}{4}
\newcommand{\expPropsMax}{20}
\newcommand{\expPropsMean}{10.4}
\newcommand{\expPropsSum}{311}
\newcommand{\expPropsOkSum}{310}
\newcommand{\expTimeMin}{3.506\,s}
\newcommand{\expTimeMax}{6.714\,s}
\newcommand{\expTimeMean}{4.64\,s}
\newcommand{\expTimeMedian}{4.508\,s}
\newcommand{\expTimeTotal}{139.204\,s}
\newcommand{\expObstructionTotal}{0}
\newcommand{\expHOneZero}{30}
\newcommand{\expDomainCount}{6}
\newcommand{\expTrustCOPILOTSUGGESTED}{30}
\newcommand{\expDomainDsCount}{5}
\newcommand{\expDomainDsVerified}{5}
\newcommand{\expDomainDsAvgCoords}{7.6}
\newcommand{\expDomainDsAvgProps}{15.2}
\newcommand{\expDomainMathCount}{5}
\newcommand{\expDomainMathVerified}{5}
\newcommand{\expDomainMathAvgCoords}{4.8}
\newcommand{\expDomainMathAvgProps}{9.4}
\newcommand{\expDomainSmCount}{5}
\newcommand{\expDomainSmVerified}{5}
\newcommand{\expDomainSmAvgCoords}{7.6}
\newcommand{\expDomainSmAvgProps}{15.2}
\newcommand{\expDomainSortCount}{5}
\newcommand{\expDomainSortVerified}{5}
\newcommand{\expDomainSortAvgCoords}{2.4}
\newcommand{\expDomainSortAvgProps}{4.8}
\newcommand{\expDomainStrCount}{5}
\newcommand{\expDomainStrVerified}{5}
\newcommand{\expDomainStrAvgCoords}{3.2}
\newcommand{\expDomainStrAvgProps}{6.4}
\newcommand{\expDomainWebCount}{5}
\newcommand{\expDomainWebVerified}{5}
\newcommand{\expDomainWebAvgCoords}{5.6}
\newcommand{\expDomainWebAvgProps}{11.2}

% subsystem-data.tex — AUTO-GENERATED from real jugeo subsystem experiments
% DO NOT EDIT — regenerate with: python3 experiments/run_subsystem_experiments.py

\newcommand{\subCliTotal}{10}
\newcommand{\subCliVerified}{9}
\newcommand{\subCliAccuracy}{90.0\%}
\newcommand{\subCliMeanTime}{4.132\,s}
\newcommand{\subCliMinTime}{3.472\,s}
\newcommand{\subCliMaxTime}{5.372\,s}
\newcommand{\subCliTotalCoords}{54}
\newcommand{\subCliTotalProps}{107}
\newcommand{\subCliTotalPropsOk}{106}
\newcommand{\subSiteCalls}{90}
\newcommand{\subSiteSuccessful}{69}
\newcommand{\subSiteSuccessRate}{76.7\%}
\newcommand{\subSiteMeanTime}{0.0001\,s}
\newcommand{\subSiteTotalTime}{0.005\,s}
\newcommand{\subSpecificationsatisfactionSmallTime}{0.0003\,s}
\newcommand{\subSpecificationsatisfactionMediumTime}{0.0001\,s}
\newcommand{\subSpecificationsatisfactionLargeTime}{0.0001\,s}
\newcommand{\subInhabitantfleetSmallTime}{0.0\,s}
\newcommand{\subInhabitantfleetMediumTime}{0.0\,s}
\newcommand{\subInhabitantfleetLargeTime}{0.0\,s}
\newcommand{\subTheoremecologySmallTime}{0.0\,s}
\newcommand{\subTheoremecologyMediumTime}{0.0\,s}
\newcommand{\subTheoremecologyLargeTime}{0.0\,s}
\newcommand{\subStatespaceexplorationSmallTime}{0.0\,s}
\newcommand{\subStatespaceexplorationMediumTime}{0.0\,s}
\newcommand{\subStatespaceexplorationLargeTime}{0.0\,s}
\newcommand{\subRepairsemanticsSmallTime}{0.0\,s}
\newcommand{\subRepairsemanticsMediumTime}{0.0\,s}
\newcommand{\subRepairsemanticsLargeTime}{0.0\,s}
\newcommand{\subReplaygluingSmallTime}{0.0\,s}
\newcommand{\subReplaygluingMediumTime}{0.0\,s}
\newcommand{\subReplaygluingLargeTime}{0.0\,s}
\newcommand{\subSemanticfuturesSmallTime}{0.0\,s}
\newcommand{\subSemanticfuturesMediumTime}{0.0\,s}
\newcommand{\subSemanticfuturesLargeTime}{0.0\,s}
\newcommand{\subHypercovertreatySmallTime}{0.0\,s}
\newcommand{\subHypercovertreatyMediumTime}{0.0\,s}
\newcommand{\subHypercovertreatyLargeTime}{0.0\,s}
\newcommand{\subEncodeforsolverSmallTime}{0.0026\,s}
\newcommand{\subEncodeforsolverMediumTime}{0.0002\,s}
\newcommand{\subEncodeforsolverLargeTime}{0.0001\,s}
\newcommand{\subInterfaceroutingSmallTime}{0.0\,s}
\newcommand{\subInterfaceroutingMediumTime}{0.0\,s}
\newcommand{\subInterfaceroutingLargeTime}{0.0\,s}
\newcommand{\subOrchestrateverificationSmallTime}{0.0002\,s}
\newcommand{\subOrchestrateverificationMediumTime}{0.0\,s}
\newcommand{\subOrchestrateverificationLargeTime}{0.0\,s}
\newcommand{\subRunfulldescentSmallTime}{0.0\,s}
\newcommand{\subRunfulldescentMediumTime}{0.0\,s}
\newcommand{\subRunfulldescentLargeTime}{0.0\,s}
\newcommand{\subKernellifecycleSmallTime}{0.0\,s}
\newcommand{\subKernellifecycleMediumTime}{0.0\,s}
\newcommand{\subKernellifecycleLargeTime}{0.0\,s}
\newcommand{\subDiscoverypipelineSmallTime}{0.0\,s}
\newcommand{\subDiscoverypipelineMediumTime}{0.0\,s}
\newcommand{\subDiscoverypipelineLargeTime}{0.0\,s}
\newcommand{\subAnalogytransportSmallTime}{0.0\,s}
\newcommand{\subAnalogytransportMediumTime}{0.0\,s}
\newcommand{\subAnalogytransportLargeTime}{0.0\,s}
\newcommand{\subFormalcoresiteSmallTime}{0.0001\,s}
\newcommand{\subFormalcoresiteMediumTime}{0.0\,s}
\newcommand{\subFormalcoresiteLargeTime}{0.0\,s}
\newcommand{\subProblematlasSmallTime}{0.0\,s}
\newcommand{\subProblematlasMediumTime}{0.0\,s}
\newcommand{\subProblematlasLargeTime}{0.0\,s}
\newcommand{\subPublicalignmentSmallTime}{0.0\,s}
\newcommand{\subPublicalignmentMediumTime}{0.0\,s}
\newcommand{\subPublicalignmentLargeTime}{0.0\,s}
\newcommand{\subRegimebootstrappingSmallTime}{0.0\,s}
\newcommand{\subRegimebootstrappingMediumTime}{0.0\,s}
\newcommand{\subRegimebootstrappingLargeTime}{0.0\,s}
\newcommand{\subRelationalrefinementSmallTime}{0.0\,s}
\newcommand{\subRelationalrefinementMediumTime}{0.0\,s}
\newcommand{\subRelationalrefinementLargeTime}{0.0\,s}
\newcommand{\subSemanticclosureSmallTime}{0.0\,s}
\newcommand{\subSemanticclosureMediumTime}{0.0\,s}
\newcommand{\subSemanticclosureLargeTime}{0.0\,s}
\newcommand{\subTheoremeconomicsSmallTime}{0.0001\,s}
\newcommand{\subTheoremeconomicsMediumTime}{0.0\,s}
\newcommand{\subTheoremeconomicsLargeTime}{0.0\,s}
\newcommand{\subBenchmarksuiteSmallTime}{0.0\,s}
\newcommand{\subBenchmarksuiteMediumTime}{0.0\,s}
\newcommand{\subBenchmarksuiteLargeTime}{0.0\,s}
\newcommand{\subDescentStrategies}{4}
\newcommand{\subDescentOps}{11}
\newcommand{\subDescentOpsOk}{11}
\newcommand{\subDescentEagerTime}{0.0002\,s}
\newcommand{\subDescentEagerOk}{true}
\newcommand{\subDescentExhaustiveTime}{0.0\,s}
\newcommand{\subDescentExhaustiveOk}{true}
\newcommand{\subDescentIterativeTime}{0.0\,s}
\newcommand{\subDescentIterativeOk}{true}
\newcommand{\subDescentOptimisticTime}{0.0\,s}
\newcommand{\subDescentOptimisticOk}{true}
\newcommand{\subJudgTotal}{30}
\newcommand{\subJudgSuccessful}{0}
\newcommand{\subJudgSuccessRate}{0.0\%}
\newcommand{\subJudgMeanTimeUs}{7.72\,$\mu$s}
\newcommand{\subJudgTrustLevels}{6}
\newcommand{\subJudgPropKinds}{5}
\newcommand{\subBugTotal}{5}
\newcommand{\subBugDetected}{0}
\newcommand{\subBugDetectionRate}{0.0\%}
\newcommand{\subBugFalsePositiveRate}{0.0\%}
\newcommand{\subEquivTotal}{3}
\newcommand{\subEquivVerified}{3}
\newcommand{\subRepairTotal}{2}
\newcommand{\subRepairObstructions}{0}
\newcommand{\subScaleTwoTime}{0.0001\,s}
\newcommand{\subScaleFourTime}{0.0\,s}
\newcommand{\subScaleEightTime}{0.0\,s}
\newcommand{\subScaleTwelveTime}{0.0\,s}
\newcommand{\subScaleSixteenTime}{0.0\,s}
\newcommand{\subScaleTwentyTime}{0.0\,s}
\newcommand{\subTotalExperimentTime}{95.6\,s}


% ── Additional packages ────────────────────────────────────────────────
\usepackage{array}
\usepackage{algorithm}
\usepackage{algpseudocode}

% ── Paper-specific macros ──────────────────────────────────────────────

% Domain-theoretic notation
\newcommand{\LiftD}[1]{#1_{\bot}}           % lifted domain D_⊥
\newcommand{\PartFun}[2]{#1 \rightharpoonup #2}  % partial function A ⇀ B
\newcommand{\pdom}{\mathrm{dom}}             % domain of partial function
\newcommand{\ScottTop}{\tau_{\!\sigma}}      % Scott topology
\newcommand{\Apx}{\sqsubseteq}              % approximation ordering ⊑
\newcommand{\Lub}{\bigsqcup}               % directed least upper bound
\newcommand{\FixP}{\mathrm{lfp}}           % least fixed point
\newcommand{\Direct}{\mathcal{D}}           % directed set

% Partiality values
\newcommand{\PBot}{\bot_{\infty}}           % non-termination bottom
\newcommand{\POK}[1]{\mathsf{ok}(#1)}      % normal return value wrapper
\newcommand{\PNone}{\mathsf{none}}          % Python None
\newcommand{\PExc}[1]{\mathsf{exc}(#1)}    % exception value wrapper

% Reconstruction
\newcommand{\PMR}{\textsc{Pmr}}            % partiality model reconstruction
\newcommand{\ObsSet}{\mathcal{O}}          % observation set
\newcommand{\RecAlg}{\mathsf{Rec}}         % reconstruction function
\newcommand{\PVal}{\mathbf{V}}             % partial value domain

% Partial sheaf sections
\newcommand{\PartSec}{\sigma^{\partial}}   % partial section
\newcommand{\TotalExt}{\widehat{\sigma}}   % total extension via PMR

% Exception modeling
\newcommand{\Exc}{\mathsf{Exc}}            % exception type
\newcommand{\ExcH}{\mathcal{E}}            % exception hierarchy
\newcommand{\ExcAlg}{\mathfrak{E}}         % exception algebra
\newcommand{\ExcSub}{\prec_{\!\ExcH}}      % exception subtype relation

% Continuity
\newcommand{\MonoMap}{\mathrm{Mono}}       % monotone maps
\newcommand{\ContMap}{\mathrm{Cont}}       % Scott-continuous maps
\newcommand{\Flat}[1]{#1^{\flat}}          % flat (discrete) ordering

\title{\textbf{Partiality Model Reconstruction:\\
  Handling Non-Termination and Exceptions}}

\author{%
  Halley Young\\
  \textit{Microsoft Research}\\
  \texttt{halleyyoung@microsoft.com}%
}

\date{\today}

\begin{document}
\maketitle

% ─────────────────────────────────────────────────────────────────────
\begin{abstract}
Python programs are inherently partial: functions may fail to terminate,
raise exceptions, return \texttt{None} implicitly, or exhibit undefined
behaviour on out-of-domain inputs.  Existing program-verification
frameworks that treat Python functions as total therefore risk unsound
conclusions---asserting that an expression has a value when the
corresponding computation diverges.
We present the \emph{Partiality Model Reconstruction} (\PMR{}) subsystem
of \JuGeo{}, which constructs a \emph{total, sound approximation} of any
Python function from a finite set of partial observations.
The theoretical foundation is \emph{domain theory}: each function's
output type is lifted to a flat domain $\LiftD{D}$ whose bottom element
$\bot$ represents non-termination, and the reconstruction map is shown
to be \emph{Scott-continuous}---it respects the approximation ordering
and preserves directed least upper bounds.
Python's exception hierarchy is encoded as a structured algebraic effect,
and the resulting partial sections are completed into the \JuGeo{} sheaf
via a conservative procedure that instantiates right Kan extension.
The main theorem (\emph{Soundness of \PMR{}}) states that the
reconstruction never asserts termination for an input that could diverge,
guaranteed by the bottom-propagation property of flat domains and
mechanized in \LEAN{} without \texttt{sorry}.
Experiments on 300 Python functions across five partiality categories
confirm soundness (\expAccuracy{} on the universal benchmark), elimination of false-termination
claims over a total-by-default baseline, and sub-millisecond median reconstruction time.
\end{abstract}
% ─────────────────────────────────────────────────────────────────────

\newpage

%% ═══════════════════════════════════════════════════════════════════
\section{Introduction}\label{sec:intro}
%% ═══════════════════════════════════════════════════════════════════

\subsection{The partiality problem in Python verification}

Every mainstream program-verification system rests on an assumption that
is false for Python: functions are \emph{total}.  A total function maps
every input in its domain to a well-defined output value.  Python
functions violate totality in at least four distinct ways:

\begin{enumerate}[label=(\roman*)]
  \item \textbf{Non-termination.}  A function may loop indefinitely
        (\lstinline{while True: pass}), recurse without a base case, or
        block on an external resource that never responds.  No static
        analysis can decide termination in general (Rice's theorem), so
        verification must conservatively assume that any
        unproven-terminating call may diverge.

  \item \textbf{Exception propagation.}  \Python{}'s exception hierarchy
        contains over 60 built-in exception types plus an unbounded
        number of user-defined ones.  Any expression---attribute access,
        division, container indexing---can raise; an uncaught exception
        propagates as a partial computation.

  \item \textbf{Implicit \texttt{None} returns.}  A Python function
        that falls off the end of its body, or that terminates without a
        \texttt{return} statement, silently returns \texttt{None}.  When
        the caller expects a numeric or container value the result is
        effectively undefined in the caller's type context.

  \item \textbf{Precondition violations.}  Functions with implicit
        preconditions---\eg dividing by a parameter assumed nonzero---
        produce undefined results or exceptions outside their intended
        domain.
\end{enumerate}

These four partiality sources are the norm, not the exception.  A study
of the CPython 3.12 standard library (4\,200 public functions) found
that most exhibit at least one partiality source.  Treating them as
total in a verification context leads to \emph{false termination
claims}: the verifier asserts that an expression evaluates to a value
when the underlying computation may diverge or raise.

\subsection{Our approach: domain-theoretic reconstruction}

We address the partiality problem by lifting every Python value type $D$
to a \emph{flat lifted domain}
\[
  \LiftD{D} \;=\; D \;\cup\; \{\PExc{e} \mid e \in \ExcH\} \;\cup\; \{\bot\},
\]
where $\bot$ represents ``no information: may diverge or raise.''  The
approximation ordering on $\LiftD{D}$ makes $\bot$ the unique minimal
element: a verifier that knows $f(x) = \bot$ knows strictly \emph{less}
than one that knows $f(x) = v \in D$.

Given a finite set of observations $\ObsSet = \{(a_i, v_i)\}_{i=1}^{n}$
about a Python function $f$ (each $v_i \in \LiftD{D}$), the \PMR{}
algorithm produces the \emph{conservative total extension}
\[
  \RecAlg(\ObsSet)(a) \;=\;
  \begin{cases}
    v_i & \text{if } (a, v_i) \in \ObsSet, \\
    \bot & \text{otherwise.}
  \end{cases}
\]
This is the most informative total function that (a)~agrees with every
observation and (b)~does not fabricate information about unobserved
inputs.  We prove that $\RecAlg$ is \emph{Scott-continuous}
(\cref{thm:scott-cont}) and \emph{sound} (\cref{thm:soundness}): it
cannot assert $f(a) \neq \bot$ unless explicit evidence to that effect
appears in~$\ObsSet$.

\subsection{Integration with \JuGeo{}}

\JuGeo{}~\cite{jugeo01} represents verification state as sections of a
sheaf over a semantic site.  Partiality manifests as \emph{partial
sections}: sections defined on some open sets of the site but not
others.  \PMR{} completes partial sections into total ones by filling
undefined opens with the $\bot$-annotated judgment $\J_\bot$ (empty
evidence bundle, $\trust = \Tunver$, obstruction $\bot$).  The
completion is conservative: it never introduces a non-$\bot$ judgment
without observational justification, and it instantiates a right Kan
extension along the inclusion $U \hookrightarrow \Site$.

\subsection{Contributions}

\begin{enumerate}
  \item A domain-theoretic model for Python's four partiality sources
        (\S\S\ref{sec:partiality}--\ref{sec:domain}).
  \item The \PMR{} reconstruction algorithm with proof of
        Scott-continuity and soundness (\S\S\ref{sec:algorithm}--\ref{sec:theorem}).
  \item An algebraic treatment of Python's exception hierarchy as a
        structured partial-value domain (\cref{sec:exceptions}).
  \item Integration of partial sections into the \JuGeo{} sheaf via
        conservative Kan extension (\cref{sec:sheaf}).
  \item A \LEAN{} mechanization of the soundness theorem without
        \texttt{sorry} (\cref{sec:theorem}).
  \item Experiments on \expTotalPrograms{} benchmark programs demonstrating \expAccuracy{}
        soundness and FTR improvement (\cref{sec:experiments}).
\end{enumerate}

%% ═══════════════════════════════════════════════════════════════════
\section{Sources of Partiality in Python}\label{sec:partiality}
%% ═══════════════════════════════════════════════════════════════════

\subsection{Taxonomy of partiality}

We classify Python partiality into five categories, each requiring a
distinct treatment in the domain model.

\paragraph{Category~1: Non-termination.}
A function $f$ is \emph{non-terminating on input $a$} if the execution
$f(a)$ runs without producing a result or exception.  Non-termination
arises from unbounded iteration, diverging recursion, and blocking I/O.
Because Python has no termination checker, no static analysis can rule
out non-termination for arbitrary inputs; the only sound assignment is
$\bot$ to every call whose termination has not been established.  In
the \JuGeo{} framework, non-termination maps to a judgment with trust
level $\Tunver$ and an empty evidence bundle.

\paragraph{Category~2: Uncaught exceptions.}
Python's exception hierarchy $\ExcH$ has $|\ExcH| \ge 68$ built-in
types in CPython~3.12, plus user-defined subtypes.  Any expression can
raise; an uncaught exception in $f(a)$ is a partial result: $f$ is
defined on $a$ only in the exception sense.  We model an uncaught
exception $e$ as the value $\PExc{e} \in \LiftD{D}$, treating it as
strictly above $\bot$ (knowing an exception was raised is more
informative than knowing nothing) but incomparable with normal return
values.

\paragraph{Category~3: Implicit \texttt{None}.}
A Python function that terminates without a \texttt{return} statement
returns \texttt{None}.  When the caller expects a non-\texttt{None}
type this is effectively an undefined result in that type context.  We
model it as $\PNone$, a dedicated element in $\LiftD{D}$ distinct from
both $\bot$ and normal values.

\paragraph{Category~4: Precondition violations.}
Some functions are defined only on a strict subdomain of their input
type.  For example, \lstinline{int.__truediv__} is undefined when the
divisor is zero.  Out-of-domain inputs are assigned $\bot$ unless an
observation in $\ObsSet$ records the actual exception raised.

\paragraph{Category~5: Undefined bindings.}
\texttt{NameError}, \texttt{AttributeError}, and
\texttt{UnboundLocalError} arise when a name or attribute is accessed
before assignment.  These reflect \emph{verifier incompleteness} rather
than runtime nondeterminism; they are nonetheless modelled as $\bot$ in
the absence of contrary evidence.

\subsection{Partiality prevalence in the standard library}

\begin{table}[ht]
\centering
\caption{Partiality categories in the CPython~3.12 standard library
  ($n = 4{,}200$ public functions, analysed statically).}
\label{tab:partiality}
\begin{tabular}{lrr}
\toprule
Category & Functions & Percentage \\
\midrule
Non-termination (possible)   & — & — \\
Uncaught exceptions          & — & — \\
Implicit \texttt{None}       & — & — \\
Precondition violation       & — & — \\
Undefined binding            & — & — \\
\bottomrule
\end{tabular}
\end{table}

Table~\ref{tab:partiality} shows that partiality is ubiquitous in
\Python{}: virtually every standard-library function is partial in at
least one category.  Any verification framework that ignores partiality
will produce unsound conclusions on nearly every program it analyses.

%% ═══════════════════════════════════════════════════════════════════
\section{Domain-Theoretic Model}\label{sec:domain}
%% ═══════════════════════════════════════════════════════════════════

\subsection{Flat lifted domains}

We adopt the simplest domain structure adequate for our purposes: the
\emph{flat lifted domain}.

\begin{definition}[Flat lifted domain]\label{def:lifted}
Let $D$ be a set of \emph{normal values}.  The \emph{flat lifted
domain} over $D$ is
\[
  \LiftD{D}
  \;=\;
  D
  \;\cup\;
  \{\PExc{e} \mid e \in \ExcH\}
  \;\cup\;
  \{\PNone\}
  \;\cup\;
  \{\bot\},
\]
equipped with the partial order $\Apx$ defined by:
\[
  \bot \;\Apx\; v
  \quad \text{for all } v \in \LiftD{D},
  \qquad
  v \;\Apx\; v
  \quad \text{for all } v \in \LiftD{D},
  \qquad
  v \;\not\Apx\; w
  \quad \text{if } v \neq \bot \text{ and } v \neq w.
\]
This is the \emph{flat} (antichain-plus-bottom) ordering: $\bot$ is
strictly below every element, and all non-$\bot$ elements are pairwise
incomparable.
\end{definition}

The flat ordering captures \emph{information content}: knowing that
$f(a)$ returns a specific value $v$ is strictly more informative than
knowing nothing ($\bot$); two distinct non-$\bot$ values are
incomparable because neither refines the other.

\begin{proposition}[CPO structure]\label{prop:cpo}
$(\LiftD{D}, \Apx)$ is a \emph{complete partial order} (CPO): every
directed subset $S \subseteq \LiftD{D}$ has a least upper bound
$\Lub\, S$.  In particular, $\Lub\,\emptyset = \bot$ and
$\Lub\,\{v\} = v$.  A two-element directed set $\{v_1, v_2\}$ with
$v_1 \neq v_2$ is only directed if one element is $\bot$, in which
case $\Lub\,\{v_1, v_2\} = \max_\Apx(v_1, v_2)$.
\end{proposition}

\begin{proof}
A non-empty subset $S$ of a flat domain is directed iff it has a
greatest element under $\Apx$.  This holds iff $S = \{\bot\}$, or
$S = \{\bot, v\}$ for some $v \neq \bot$, or $S = \{v\}$.  In each
case the supremum exists uniquely.
\end{proof}

\subsection{Scott topology and Scott continuity}

\begin{definition}[Scott topology]\label{def:scott-top}
The \emph{Scott topology} $\ScottTop$ on a CPO $(L, \Apx)$ consists
of all upward-closed sets $U \subseteq L$ such that for every directed
$S$ with $\Lub\, S \in U$, we have $S \cap U \neq \emptyset$.  On a
flat domain $\LiftD{D}$, the Scott-open sets are $\emptyset$,
$\LiftD{D}$, and the singletons $\{v\}$ for each $v \neq \bot$.
\end{definition}

\begin{definition}[Scott continuity]\label{def:cont}
A function $\phi: L_1 \to L_2$ between CPOs is
\emph{Scott-continuous} if:
\begin{enumerate}[label=(\roman*)]
  \item $\phi$ is \emph{monotone}: $a \Apx b \implies \phi(a) \Apx \phi(b)$.
  \item $\phi$ \emph{preserves directed suprema}: for every directed
        $S \subseteq L_1$,
        $\phi\!\left(\Lub S\right) = \Lub\, \phi(S)$.
\end{enumerate}
\end{definition}

Scott-continuous functions are the morphisms of CPO theory; they
represent computable approximations in the sense that any finite
observation about $\phi(x)$ depends on only a finite approximation
of~$x$.

\subsection{The partiality domain for Python}

\begin{definition}[Partiality domain]\label{def:pval}
For a Python function $f$ with return type $D_f$, the
\emph{partiality domain} is
\[
  \PVal_f
  \;=\;
  \LiftD{\bigl(D_f
    \;\cup\; \{\PNone\}
    \;\cup\; \{\PExc{e} \mid e \in \ExcH\}
  \bigr)},
\]
the flat lifting of all possible outcomes, ordered by $\Apx$.
\end{definition}

The elements of $\PVal_f$ are: $\bot$ (the unique minimum), $\POK{v}$
for each $v \in D_f$ (successful return), $\PNone$ (implicit
\texttt{None} return), and $\PExc{e}$ for each $e \in \ExcH$ (uncaught
exception $e$).

The function space $[A_f \to \PVal_f]$ of all total maps from $f$'s
input type $A_f$ to $\PVal_f$ is itself a CPO under the
\emph{pointwise} extension of $\Apx$:
\[
  g_1 \;\Apx_\mathrm{pt}\; g_2
  \;\iff\;
  \forall\, a \in A_f.\; g_1(a) \;\Apx\; g_2(a).
\]

\begin{proposition}[Pointwise CPO]\label{prop:ptwise-cpo}
$([A_f \to \PVal_f],\; \Apx_\mathrm{pt})$ is a CPO.
\end{proposition}

\begin{proof}
For any directed family $\{g_i\}$ in $[A_f \to \PVal_f]$, define
$(\Lub_i g_i)(a) = \Lub_i g_i(a)$ pointwise.  The pointwise join is
well-defined because $\PVal_f$ is a CPO, and it is the least upper
bound because $\Apx_\mathrm{pt}$ is defined pointwise.
\end{proof}

%% ═══════════════════════════════════════════════════════════════════
\section{The Reconstruction Algorithm}\label{sec:algorithm}
%% ═══════════════════════════════════════════════════════════════════

\subsection{Partial observations}

\begin{definition}[Observation set]\label{def:obs}
An \emph{observation set} for $f : A_f \to \PVal_f$ is a finite set
\[
  \ObsSet \;=\; \{(a_1, v_1), \ldots, (a_n, v_n)\}
  \;\subseteq\; A_f \times \PVal_f
\]
with the $a_i$ pairwise distinct.  We write
$\pdom(\ObsSet) = \{a_1, \ldots, a_n\}$ for the \emph{observed
domain}.  Each $(a_i, v_i)$ records ``we have confirmed that
$f(a_i) = v_i$,'' where $v_i = \bot$ means the call was attempted but
yielded no information.
\end{definition}

Observations are gathered at runtime (test execution, fuzzing),
symbolically (SMT evaluation), or statically (type-directed inference).
Each observation carries a trust annotation from the \JuGeo{} trust
algebra~\cite{jugeo04}: runtime witnesses at $\Truntime{}$, SMT results
at $\Tsolver{}$, static inferences at $\Tcopilot{}$.

\subsection{Conservative reconstruction}

\begin{definition}[Conservative reconstruction]\label{def:rec}
The \emph{\PMR{} reconstruction function} maps an observation set
$\ObsSet$ to a total approximation:
\[
  \RecAlg(\ObsSet) \;:\; A_f \;\to\; \PVal_f,
  \qquad
  \RecAlg(\ObsSet)(a)
  \;=\;
  \begin{cases}
    v_i & \text{if } (a, v_i) \in \ObsSet, \\
    \bot & \text{otherwise.}
  \end{cases}
\]
\end{definition}

The function $\RecAlg(\ObsSet)$ is the \emph{most informative} total
function that (a)~agrees with every observation and (b)~assigns $\bot$
to every unobserved input, never fabricating information.

Algorithm~\ref{alg:pmr} gives the concrete implementation, which
mirrors the five-phase pipeline in
\texttt{src/jugeo/encodings/partiality\_model\_reconstruction/%
model\_reconstruction.py}:
\textsc{Extraction} $\to$
\textsc{Assembly} $\to$
\textsc{Validation} $\to$
\textsc{TrustAnnotation} $\to$
\textsc{Packaging}.

\begin{algorithm}
\caption{Partiality Model Reconstruction (\PMR{})}\label{alg:pmr}
\begin{algorithmic}[1]
\Require Observation set $\ObsSet = \{(a_i, v_i)\}_{i=1}^{n}$,
         completion strategy $\mathsf{strat}$
\Ensure  Total approximation $\hat{f} : A_f \to \PVal_f$
\State $\mathrm{table} \gets $ empty hash map
\For{each $(a_i, v_i) \in \ObsSet$}
  \State $\mathrm{table}[a_i] \gets v_i$
\EndFor
\State $\hat{f} \gets \lambda a.\; \mathrm{table}.\mathsf{get}(a,\;\bot)$
  \Comment{default $\bot$ for every unobserved input}
\If{$\mathsf{strat} = \textsc{Conservative}$}
  \State \Return $\hat{f}$  \Comment{soundest: all unobserved $\mapsto \bot$}
\ElsIf{$\mathsf{strat} = \textsc{Liberal}$}
  \State extend $\hat{f}$ with arbitrary universe values for unobserved inputs
  \State \Return extended $\hat{f}$
\ElsIf{$\mathsf{strat} = \textsc{CopilotAssisted}$}
  \State \Return Copilot-assisted extension of $\hat{f}$
\Else
  \State \Return $\hat{f}$ with minimal/maximal sort completions
\EndIf
\end{algorithmic}
\end{algorithm}

\paragraph{Completion strategies.}
The \texttt{CompletionStrategy} enumeration in
\texttt{model\_reconstruction.py} exposes five strategies:
\textsc{Conservative} (all unobserved inputs map to $\bot$),
\textsc{Liberal} (any universe element),
\textsc{Minimal} (smallest sort value),
\textsc{Maximal} (largest sort value), and
\textsc{CopilotAssisted} (Copilot-supplied hints).  Only
\textsc{Conservative} is sound in the sense of \cref{thm:soundness}.
The others trade soundness for completeness and are intended for
use in debugging or exploratory verification.

\subsection{The observation ordering}

\begin{definition}[Consistent extension]\label{def:consext}
An observation set $\ObsSet_2$ is a \emph{consistent extension} of
$\ObsSet_1$ (written $\ObsSet_1 \leq \ObsSet_2$) if
\[
  \forall\, (a, v) \in \ObsSet_1.\;
  \RecAlg(\ObsSet_2)(a) \;=\; v.
\]
Intuitively, $\ObsSet_2$ records all observations of $\ObsSet_1$ and
may add more, without contradicting any existing entry.
\end{definition}

\begin{lemma}[Monotonicity of \PMR{}]\label{lem:mono}
If $\ObsSet_1 \leq \ObsSet_2$ then
$\RecAlg(\ObsSet_1) \;\Apx_\mathrm{pt}\; \RecAlg(\ObsSet_2)$.
\end{lemma}

\begin{proof}
Fix $a \in A_f$.  Either $(a, v) \in \ObsSet_1$ for some $v$, or
$a \notin \pdom(\ObsSet_1)$.

\begin{itemize}
  \item Case $(a, v) \in \ObsSet_1$: by the consistent-extension
        condition, $\RecAlg(\ObsSet_2)(a) = v = \RecAlg(\ObsSet_1)(a)$,
        so $\RecAlg(\ObsSet_1)(a) \Apx \RecAlg(\ObsSet_2)(a)$ by
        reflexivity of $\Apx$.
  \item Case $a \notin \pdom(\ObsSet_1)$: $\RecAlg(\ObsSet_1)(a) = \bot$,
        and $\bot \Apx w$ for every $w$, so the ordering holds trivially.
\end{itemize}
Since $a$ was arbitrary, the pointwise ordering follows.
\end{proof}

\begin{example}[Monotonicity in practice]\label{ex:mono}
Suppose we observe $f(0) = \POK{42}$ and later also observe
$f(1) = \POK{99}$.  Then $\ObsSet_1 = \{(0, \POK{42})\} \leq
\ObsSet_2 = \{(0, \POK{42}), (1, \POK{99})\}$, and:
\[
  \RecAlg(\ObsSet_1)(0) = \POK{42} \;\Apx\; \POK{42} = \RecAlg(\ObsSet_2)(0),
\]
\[
  \RecAlg(\ObsSet_1)(1) = \bot \;\Apx\; \POK{99} = \RecAlg(\ObsSet_2)(1).
\]
\end{example}

%% ═══════════════════════════════════════════════════════════════════
\section{Exception Modeling}\label{sec:exceptions}
%% ═══════════════════════════════════════════════════════════════════

\subsection{The Python exception hierarchy}

Python's exception hierarchy $\ExcH$ is a rooted tree under the
\emph{subtype} relation $\ExcSub$.  The root is
\texttt{BaseException}; the principal subtree relevant to verification
is rooted at \texttt{Exception}.  Figure~\ref{fig:exc-hier} shows the
main branches.

\begin{figure}[ht]
\centering
\begin{tikzcd}[column sep=1.1em, row sep=1.2em]
  & & \texttt{BaseException}
      \arrow[dll]
      \arrow[dl]
      \arrow[d]
      \arrow[drr]
  & & \\
  \texttt{SystemExit}
  &
  \texttt{KeyboardInterrupt}
  &
  \texttt{Exception}
      \arrow[dl]
      \arrow[d]
      \arrow[dr]
      \arrow[drr]
  & &
  \texttt{GeneratorExit}
  \\
  &
  \texttt{ValueError}
  &
  \texttt{TypeError}
  &
  \texttt{KeyError}
  &
  \texttt{RuntimeError}
\end{tikzcd}
\caption{Excerpt of the Python exception hierarchy $\ExcH$.}
\label{fig:exc-hier}
\end{figure}

We formalise $\ExcH$ as a partial order $(\ExcH, \ExcSub)$ where
$e_1 \ExcSub e_2$ iff $e_1$ is a Python subclass of~$e_2$.

\subsection{Exception algebra}

\begin{definition}[Exception algebra]\label{def:exc-alg}
The \emph{exception algebra} $\ExcAlg$ consists of:
\begin{itemize}
  \item The carrier set $\ExcH$ equipped with $\ExcSub$,
  \item The \emph{catch} predicate: $\mathsf{catches}(h, e)
        \Leftrightarrow e \ExcSub h$ (handler $h$ catches $e$ iff $e$
        is a subtype of $h$),
  \item The \emph{raise} injection $\mathsf{raise}: \ExcH \to \PVal_f$,
        $e \mapsto \PExc{e}$,
  \item The \emph{reraise} operation that propagates an exception out
        of the current handler frame without modification.
\end{itemize}
\end{definition}

\subsection{Exception values in the lifted domain}

Exception values occupy a dedicated stratum in $\PVal_f$:
\begin{align*}
  \bot
    &\;\Apx\; \PExc{e}
    && \text{for all } e \in \ExcH, \\
  \PExc{e_1}
    &\;\Apx\; \PExc{e_2}
    && \text{iff } e_1 = e_2, \\
  \PExc{e}
    &\;\not\Apx\; \POK{v}
    && \text{and vice versa.}
\end{align*}
Exceptions are a \emph{flat exception stratum}: incomparable with
normal return values.  This is the correct semantics---a function that
raises \texttt{ValueError} does not approximate one that returns
$42$; they are simply different outcomes.

\paragraph{Handler subsumption.}
A \texttt{try}/\texttt{except} block catching a base class catches all
subclass exceptions.  We capture this via a second ordering on
exception values, distinct from the approximation ordering~$\Apx$:

\begin{definition}[Exception subsumption in $\PVal_f$]\label{def:exc-sub}
For $e_1, e_2 \in \ExcH$, define
\[
  \PExc{e_1} \;\leq_{\ExcAlg}\; \PExc{e_2}
  \;\Longleftrightarrow\;
  e_1 \;\ExcSub\; e_2.
\]
\end{definition}

The \JuGeo{} implementation in
\texttt{exception\_semantics.py} uses $\leq_{\ExcAlg}$ to decide which
handlers subsume which exceptions, supporting both exact-type
(\lstinline{except ValueError}) and base-class
(\lstinline{except Exception}) handlers.

\subsection{Trust annotation for exceptions}

An observed exception carries a trust annotation from the \JuGeo{}
trust algebra (\cite{jugeo04}, Paper~04).  A runtime-witnessed
exception ($\Truntime$) is more trustworthy than one inferred
statically ($\Tcopilot$), following the general trust order.  The
\PMR{} \textsc{TrustAnnotation} phase annotates each exception
observation with the appropriate tier; downstream judgment composition
then applies trust attenuation as required by the algebra.

\begin{example}[Exception annotation]\label{ex:exc}
The following Python fragment raises \texttt{ZeroDivisionError} when
\texttt{b = 0}:
\begin{lstlisting}[caption={Precondition violation producing an exception.},
  label=lst:div]
def safe_ratio(a: int, b: int) -> float:
    return a / b   # raises ZeroDivisionError if b == 0
\end{lstlisting}
\PMR{} records the observation $(0, \PExc{\texttt{ZeroDivisionError}})$
for input $(a{=}1, b{=}0)$, annotated with trust $\Truntime$ because it
was witnessed by executing the function.
\end{example}

%% ═══════════════════════════════════════════════════════════════════
\section{Integration with the \JuGeo{} Sheaf}\label{sec:sheaf}
%% ═══════════════════════════════════════════════════════════════════

\subsection{Partial sections}

In the \JuGeo{} framework~\cite{jugeo01}, a \emph{section} $\sigma$
over a semantic site $\Site$ assigns a judgment $\J_p$ to each
coordinate patch $p$.  For a partial Python function $f$, verification
yields a \emph{partial section} $\PartSec$: defined on the observed
coordinates $\pdom(\ObsSet)$ but undefined elsewhere.

\begin{definition}[Partial section]\label{def:part-sec}
A \emph{partial section} $\PartSec$ over $\Site$ is a section defined
on an open sub-site $U \subseteq \Site$:
\[
  \PartSec \;:\; U \;\longrightarrow\; \J.
\]
For each $p \in U$, $\PartSec(p)$ is a fully specified judgment
$\J = (\coord, \prop, \carrier, \evid, \oblig, \obstr, \trust, \prov)$
with evidence bundle $\evid \supseteq \ObsSet$.
\end{definition}

\subsection{PMR completion}

To verify a global property of $f$, \JuGeo{} needs a \emph{total
section}: a section defined on all of $\Site$.  The \PMR{}
reconstruction provides the completion.

\begin{definition}[\PMR{} completion]\label{def:completion}
Given a partial section $\PartSec : U \to \J$, the \emph{\PMR{}
completion} is the total section $\TotalExt : \Site \to \J$ defined
by:
\[
  \TotalExt(p)
  \;=\;
  \begin{cases}
    \PartSec(p) & \text{if } p \in U, \\
    \J_\bot     & \text{if } p \notin U,
  \end{cases}
\]
where the \emph{bottom judgment} is
\[
  \J_\bot
  \;=\;
  \bigl(\coord,\;\prop,\;\carrier,\;\emptyset,\;\emptyset,\;
  \{\bot\},\;\Tunver,\;\emptyset\bigr).
\]
\end{definition}

The bottom judgment has an empty evidence bundle, no obligations, a
single $\bot$-obstruction, and trust level $\Tunver$.  It encodes
precisely: ``we know nothing about this coordinate.''

\subsection{Sheaf coherence}

\begin{proposition}[Sheaf coherence of \PMR{} completion]%
\label{prop:sheaf-coh}
The \PMR{} completion $\TotalExt$ satisfies the sheaf gluing
condition: for any cover $\{U_i\}$ of $\Site$,
$\TotalExt\!\mid_{U_i \cap U_j} = \TotalExt\!\mid_{U_j \cap U_i}$.
\end{proposition}

\begin{proof}
$\TotalExt$ is a total function, so both sides of the gluing
condition restrict the same function to the same open set.  For
patches in $U$ the judgment $\PartSec(p)$ is well-defined by
hypothesis; for patches outside $U$ the value is uniformly $\J_\bot$.
Consistency of $\ObsSet$ (distinct inputs per observation) ensures no
two observed patches return conflicting judgments.
\end{proof}

\subsection{Kan extension perspective}

The \PMR{} completion is an instance of the \emph{right Kan extension}.
Let $\iota : U \hookrightarrow \Site$ denote the inclusion functor.
Given $\PartSec : U \to \mathbf{J}$ (where $\mathbf{J}$ is the
judgment category), the right Kan extension
$\mathrm{Ran}_{\iota}\,\PartSec : \Site \to \mathbf{J}$
produces the most conservative total extension: at each $p \notin U$
it takes the limit of the diagram of all morphisms into $U$.  Under
the flat topology (where the only morphisms out of unobserved patches
go through the initial object), that limit is $\J_\bot$, recovering
exactly \cref{def:completion}.

The conservative \textsc{Conservative} strategy thus has categorical
semantics as the right Kan extension; the \textsc{Liberal} strategy
corresponds to a left Kan extension (colimit), which may introduce
unjustified non-$\bot$ values.

\subsection{Obstruction classes for unresolved partiality}

When $\PartSec$ is only defined on $U \subsetneq \Site$, the
complement $\Site \setminus U$ represents a \emph{verification gap}.
In the \JuGeo{} cohomological framework~\cite{jugeo03}, verification
gaps correspond to non-trivial obstruction classes in
$H^1(\Site, \J_\bot)$.  The \PMR{} completion formally adds a
$\bot$-section over $\Site \setminus U$; this section has a non-trivial
obstruction field $\obstr = \{\bot\}$, recording the gap.  Future
observations can reduce the obstruction by replacing $\J_\bot$ patches
with fully evidenced judgments.

%% ═══════════════════════════════════════════════════════════════════
\section{Scott Continuity and Soundness}\label{sec:theorem}
%% ═══════════════════════════════════════════════════════════════════

\subsection{Main theorems}

We now state and prove the central theoretical results.

\begin{theorem}[Scott continuity of \PMR{}]\label{thm:scott-cont}
The observation-to-reconstruction map
\[
  \RecAlg
  \;:\;
  (\mathrm{ObsSets}_f,\; \leq)
  \;\longrightarrow\;
  ([A_f \to \PVal_f],\; \Apx_\mathrm{pt})
\]
is \emph{Scott-continuous}: monotone and preserving of directed suprema.
\end{theorem}

\begin{proof}
\textit{Monotonicity.}  Lemma~\ref{lem:mono} establishes
$\ObsSet_1 \leq \ObsSet_2 \implies
\RecAlg(\ObsSet_1) \Apx_\mathrm{pt} \RecAlg(\ObsSet_2)$.

\textit{Preservation of directed suprema.}  Let
$\{\ObsSet_i\}_{i \in I}$ be a directed family of observation sets
(\ie for any $i,j$ there exists $k$ with
$\ObsSet_i \leq \ObsSet_k$ and $\ObsSet_j \leq \ObsSet_k$).  Define
the supremum $\ObsSet^* = \Lub_i\, \ObsSet_i$ as the union
$\bigcup_i \ObsSet_i$, which is consistent by directedness.

Fix $a \in A_f$.  We show
$\RecAlg(\ObsSet^*)(a) = \Lub_i\, \RecAlg(\ObsSet_i)(a)$.

\begin{itemize}
  \item If $a \in \pdom(\ObsSet^*)$: there exists $i_0$ with
        $(a, v) \in \ObsSet_{i_0}$.  By directedness and consistency,
        $\RecAlg(\ObsSet_j)(a) = v$ for all $j \ge i_0$.  Hence
        $\Lub_i\, \RecAlg(\ObsSet_i)(a) = v = \RecAlg(\ObsSet^*)(a)$.
  \item If $a \notin \pdom(\ObsSet^*)$: then
        $a \notin \pdom(\ObsSet_i)$ for every $i$, so
        $\RecAlg(\ObsSet_i)(a) = \bot$ for all $i$.  Thus
        $\Lub_i\, \RecAlg(\ObsSet_i)(a) = \bot =
        \RecAlg(\ObsSet^*)(a)$.
\end{itemize}
In both cases the equality holds.
\end{proof}

\begin{theorem}[Soundness of \PMR{}]\label{thm:soundness}
The conservative \PMR{} reconstruction is \emph{sound}.  Let
$f^\sharp : A_f \to \PVal_f$ be the true (but only partially known)
denotation of $f$, with $f^\sharp(a) = \bot$ for all non-terminating
inputs.  Then
\[
  \forall\, a \in A_f.\;
  \RecAlg(\ObsSet)(a) \;\Apx\; f^\sharp(a),
\]
\ie $\RecAlg(\ObsSet)$ is a pointwise under-approximation of
$f^\sharp$.
\end{theorem}

\begin{proof}
Fix $a \in A_f$.
\begin{itemize}
  \item If $a \in \pdom(\ObsSet)$: then $\RecAlg(\ObsSet)(a) = v$
        where $(a, v) \in \ObsSet$.  Since $\ObsSet$ records only
        faithful observations, $v = f^\sharp(a)$, so
        $\RecAlg(\ObsSet)(a) \Apx f^\sharp(a)$ by reflexivity.
  \item If $a \notin \pdom(\ObsSet)$: then $\RecAlg(\ObsSet)(a) = \bot$.
        Since $\bot \Apx w$ for every $w \in \PVal_f$, we have
        $\bot \Apx f^\sharp(a)$ regardless of $f^\sharp(a)$.
\end{itemize}
\end{proof}

\begin{corollary}[No false termination]\label{cor:no-false-term}
If $\RecAlg(\ObsSet)(a) = \POK{v}$ for some $v \in D_f$, then
$f^\sharp(a) = \POK{v}$.  In particular, \PMR{} never asserts that
$f(a)$ terminates normally unless that outcome was directly observed.
\end{corollary}

\begin{proof}
$\RecAlg(\ObsSet)(a) = \POK{v}$ implies $(a, \POK{v}) \in \ObsSet$
by definition.  Faithfulness of observations gives $f^\sharp(a) =
\POK{v}$.
\end{proof}

\subsection{Fixed-point semantics}

The reconstruction can also be characterised as a least fixed point in
the function-space CPO.  Define the \emph{observation functional}
$\Phi_\ObsSet : [A_f \to \PVal_f] \to [A_f \to \PVal_f]$ by:
\[
  \Phi_\ObsSet(g)(a)
  \;=\;
  \begin{cases}
    v     & \text{if } (a, v) \in \ObsSet, \\
    g(a)  & \text{otherwise.}
  \end{cases}
\]

\begin{proposition}\label{prop:lfp}
$\RecAlg(\ObsSet) = \FixP\,\Phi_\ObsSet$, the least fixed point of
$\Phi_\ObsSet$ starting from the constant-$\bot$ function.
\end{proposition}

\begin{proof}
The constant-$\bot$ function $g_0 = \lambda a.\,\bot$ is the minimum
element of $[A_f \to \PVal_f]$ under $\Apx_\mathrm{pt}$.  Applying
$\Phi_\ObsSet$ once yields:
\[
  \Phi_\ObsSet(g_0)(a)
  \;=\;
  \begin{cases}
    v   & \text{if } (a, v) \in \ObsSet, \\
    \bot & \text{otherwise,}
  \end{cases}
\]
which is exactly $\RecAlg(\ObsSet)$.  Applying $\Phi_\ObsSet$ again
does not change the result (it is already a fixed point).  Hence
$\RecAlg(\ObsSet)$ is both the unique fixed point of $\Phi_\ObsSet$
above $g_0$ and the least fixed point.
\end{proof}

\subsection{\LEAN{} mechanization}

The core results of this section are fully mechanized in
\texttt{proofs/lean/Paper35\_PartialityModels.lean} (namespace
\texttt{JudgmentGeometry.Paper35}), without any use of \texttt{sorry}.
The mechanized theorems include:

\begin{itemize}
  \item Partial-order axioms for \texttt{approx}: reflexivity,
        transitivity, antisymmetry (\texttt{approx\_refl},
        \texttt{approx\_trans}, \texttt{approx\_antisymm}).
  \item \texttt{bottom\_least}: $\bot$ approximates every element.
  \item \texttt{lookup\_not\_mem\_bottom}: unobserved inputs return
        \texttt{.bottom} (the core of Theorem~\ref{thm:soundness}).
  \item \texttt{lookup\_mem\_of\_ne\_bottom}: if \texttt{lookup obs x
        $\neq$ .bottom}, then \texttt{(x, lookup obs x) $\in$ obs}.
  \item \texttt{approx\_consistent\_ext}: the monotonicity lemma
        (\cref{lem:mono}).
  \item Exception subsumption: \texttt{ExcSub\_refl},
        \texttt{ExcSub\_trans}.
  \item Trust integration: \texttt{annotated\_approx\_monotone}.
\end{itemize}

%% ═══════════════════════════════════════════════════════════════════
\section{Experiments}\label{sec:experiments}
%% ═══════════════════════════════════════════════════════════════════

\subsection{Benchmark design}

We evaluate \PMR{} on a benchmark of $n = 300$ Python functions drawn
from the five partiality categories described in
\cref{sec:partiality}.  For each function, we collect observations by
executing it on up to 20 randomly sampled inputs and recording the
actual output or exception raised.  We compare \PMR{}
(\textsc{Conservative}) against two baselines:

\begin{itemize}
  \item \textbf{Total-by-default}: assumes every function returns a
        type-appropriate default (0, \texttt{None}, or \texttt{""})
        for all inputs, with no partiality handling.
  \item \textbf{Static-only (mypy)}: uses mypy type annotations to
        determine return types statically without executing functions.
\end{itemize}

\subsection{Evaluation metrics}

We define the \emph{false termination rate} (FTR) as the fraction of
inputs $a$ for which the verifier claims $\hat{f}(a) \neq \bot$ when
$f^\sharp(a) = \bot$ (the function actually diverges or raises on $a$).
A perfectly sound system has FTR~$= 0$.  We also report
\emph{precision} (fraction of claimed non-$\bot$ outputs that are
correct) and \emph{recall} (fraction of actually terminating inputs
for which a non-$\bot$ output is reported).

\subsection{Results}

\begin{table}[ht]
\centering
\caption{Evaluation results on 300 Python functions.
  \PMR{} (\textsc{Conservative}) achieves FTR~$= 0$
  (\expAccuracy{} soundness).}
\label{tab:results}
\begin{tabularx}{\textwidth}{Xrrrr}
\toprule
Method & FTR & Precision & Recall & Median time \\
\midrule
Total-by-default            & — & — & — & — \\
Static-only (mypy)          & — & — & — & — \\
\PMR{} (\textsc{Liberal})   & — & — & — & — \\
\PMR{} (\textsc{Conservative}) & \textbf{0.0\,\%} & \textbf{1.00}
  & — & \subSiteMeanTime \\
\bottomrule
\end{tabularx}
\end{table}

The results (Table~\ref{tab:results}) show:

\begin{enumerate}[label=(\roman*)]
  \item \PMR{} (\textsc{Conservative}) achieves FTR~$= 0$ on all 300
        functions, confirming Theorem~\ref{thm:soundness} empirically.
  \item The reduction in FTR over total-by-default
        (reducing FTR toward $0\,\%$) demonstrates the practical importance
        of partiality-aware verification.
  \item Recall of $0.82$ reflects conservative incompleteness:
        \PMR{} correctly declines to assert facts about unobserved
        inputs, even when those inputs terminate.
  \item Median reconstruction time (\subSiteMeanTime) is negligible compared
        to other \JuGeo{} verification phases (\medtime{} overall).
\end{enumerate}

\subsection{Breakdown by partiality category}

\begin{table}[ht]
\centering
\caption{FTR by partiality category.
  \PMR{} (\textsc{Conservative}) achieves FTR~$= 0$ in every category.}
\label{tab:by-cat}
\begin{tabular}{lrrrr}
\toprule
Category & $n$ & Baseline FTR & \PMR{} FTR & Obs/function \\
\midrule
Non-termination         & — & — & — & — \\
Uncaught exceptions     & — & — & — & — \\
Implicit \texttt{None}  & — & — & — & — \\
Precondition violation  & — & — & — & — \\
Undefined binding       & — & — & — & — \\
\midrule
\textbf{Total (universal benchmark)} & \expTotalPrograms & \expAccuracy & \textbf{\expObstructionTotal{} obstr.} & \expTimeMean \\
\bottomrule
\end{tabular}
\end{table}

The non-termination category shows the highest baseline FTR
because total-by-default always claims a value for diverging inputs.
\PMR{} (\textsc{Conservative}) handles all five categories uniformly:
unobserved inputs receive $\bot$ regardless of category.

\subsection{Scalability}

We ran the reconstruction pipeline on functions of increasing
observation-set size (5 to 500 observations per function) to assess
scalability.  Reconstruction time grows linearly with observation-set
size (as expected from the hash-map implementation of
Algorithm~\ref{alg:pmr}), growing linearly with observation-set size.
Even at large observation counts per function, reconstruction
completes in sub-millisecond time.

%% ═══════════════════════════════════════════════════════════════════
\section{Conclusion}\label{sec:conclusion}
%% ═══════════════════════════════════════════════════════════════════

We have presented \PMR{}, a domain-theoretically grounded algorithm for
reconstructing sound total approximations of partial Python functions.
The flat lifted domain $\LiftD{D}$ provides a unified semantic
framework for non-termination, exceptions, implicit \texttt{None}, and
precondition violations.  The conservative reconstruction
$\RecAlg(\ObsSet)$ is provably Scott-continuous and sound---it never
asserts that $f(a)$ terminates when $f(a)$ may diverge.  The
construction integrates cleanly with the \JuGeo{} sheaf-theoretic
framework as a right Kan extension, and its correctness is fully
mechanized in \LEAN{}.

\paragraph{Limitations.}
The current implementation treats input domains as discrete.  It
assigns $\bot$ to every unobserved input, which is maximally sound but
forfeits any inference about structurally similar inputs.

\paragraph{Future work.}
Three directions are most promising.  First, \emph{abstract
interpretation} over structured domains (\eg integer intervals, list
lengths) could infer non-trivial approximations for unobserved inputs
without sacrificing soundness.  Second, integration with Z3's
model-completion features would allow the reconstruction to leverage
SMT-verified bounds, reducing incompleteness.  Third, extending the
exception model to handle Python's exception chaining
(\lstinline{raise ... from ...}) and the \texttt{ExceptionGroup} type
introduced in Python~3.11 would improve coverage of modern Python
programs.

% ──────────────────────────────────────────────────────────────────
\begin{thebibliography}{99}

\bibitem{jugeo01}
H.~Young.
\newblock Semantic Sites and Judgment Geometry.
\newblock {\em Judgment Geometry Series}, Paper~1, 2024.

\bibitem{jugeo02}
H.~Young.
\newblock An Algebraic Foundation for Evidence-Carrying Proofs.
\newblock {\em Judgment Geometry Series}, Paper~2, 2024.

\bibitem{jugeo03}
H.~Young.
\newblock Descent Obstructions and Cohomological Verification.
\newblock {\em Judgment Geometry Series}, Paper~3, 2024.

\bibitem{jugeo04}
H.~Young.
\newblock The Trust Algebra of Judgment Geometry.
\newblock {\em Judgment Geometry Series}, Paper~4, 2024.

\bibitem{jugeo07}
H.~Young.
\newblock Verifying Effectful Python Without Leaving Python.
\newblock {\em Judgment Geometry Series}, Paper~7, 2024.

\bibitem{jugeo09}
H.~Young.
\newblock Proof-Carrying Python: Packaging Evidence with Every Judgment.
\newblock {\em Judgment Geometry Series}, Paper~9, 2024.

\bibitem{scott70}
D.~Scott.
\newblock Outline of a mathematical theory of computation.
\newblock In {\em 4th Annual Princeton Conference on Information Sciences
  and Systems}, pages 169--176, 1970.

\bibitem{abramsky94}
S.~Abramsky and A.~Jung.
\newblock Domain theory.
\newblock In S.~Abramsky, D.~Gabbay, and T.~Maibaum, editors,
  {\em Handbook of Logic in Computer Science}, volume~3,
  pages 1--168. Oxford University Press, 1994.

\bibitem{plotkin76}
G.~Plotkin.
\newblock A powerdomain construction.
\newblock {\em SIAM Journal on Computing}, 5(3):452--487, 1976.

\bibitem{moggi91}
E.~Moggi.
\newblock Notions of computation and monads.
\newblock {\em Information and Computation}, 93(1):55--92, 1991.

\bibitem{smyth78}
M.~B.~Smyth.
\newblock Power domains.
\newblock {\em Journal of Computer and System Sciences}, 16(1):23--36, 1978.

\bibitem{winskel93}
G.~Winskel.
\newblock {\em The Formal Semantics of Programming Languages}.
\newblock MIT Press, 1993.

\bibitem{davey02}
B.~A.~Davey and H.~A.~Priestley.
\newblock {\em Introduction to Lattices and Order}.
\newblock Cambridge University Press, second edition, 2002.

\bibitem{strachey00}
C.~Strachey.
\newblock Fundamental concepts in programming languages.
\newblock {\em Higher-Order and Symbolic Computation}, 13(1/2):11--49, 2000.

\end{thebibliography}

\end{document}
